\documentclass{book}
\usepackage{ctex}
\usepackage[cmex10]{amsmath}
\usepackage{xcolor}
\usepackage{mathtools}
\usepackage{booktabs}
\DeclareRobustCommand{\bluebold}[1]{\textbf{\textcolor[rgb]{0.31,0.63,0.86}{#1}}}

\usepackage[
    pdfpagemode=UseOutlines,
    bookmarks=true,
    bookmarksopen=true,
    bookmarksnumbered=true,
    colorlinks=true,
    linkcolor=black,
    citecolor=black,
    urlcolor=black
]{hyperref}
\usepackage{bookmark}

\raggedbottom

\begin{document}
\tableofcontents
\clearpage

流体在管内的流动阻力

\begin{table}[htbp]
\centering
\begin{tabular}{c c} 
\toprule
知识点 & point\\ [0.5ex] 
\midrule
1 & 直管阻力损失 \\
2 & 局部阻力损失 \\
3 & 总阻力损失的计算 \\
\bottomrule
\end{tabular}
\caption{Table of Knowledge Points on Fluid Resistance Inside Pipelines.}
\label{table:data}
\end{table}

\chapter{直管阻力损失}

实际流体具有粘性，这是导致流体流动产生阻力损失的\bluebold{根本原因}。

\begin{equation}
h_\mathrm{f} = \lambda \frac{l}{d} \frac{u^2}{2}
\label{eq:fanning}
\end{equation}

Fanning 公式或阻力计算通式，适用层流、湍流。

$\lambda$：摩擦系数（因数），量纲为1，没有单位；


层流流动
\begin{equation}
\lambda = \frac{64}{\mathrm{Re}}
\label{eq:lambda-laminar}
\end{equation}

湍流流动
\begin{equation}
\lambda = f\left(\frac{\varepsilon}{d},\; \mathrm{Re}\right)
\label{eq:lambda-turbulent}
\end{equation}

$\varepsilon$：管壁绝对粗糙度，m;

$\frac{\varepsilon}{d}$：相对粗糙度

湍流区：$\lambda = f\left(\frac{\varepsilon}{d},\; \mathrm{Re}\right),\; \mathrm{Re} up, \lambda down$

阻力平方区：$\lambda = f(\frac{\varepsilon}{d})$

$\lambda$计算：查图或公式计算，图或公式一般会给

\chapter{局部阻力损失}

实际露天具有粘性，这是导致流体流体产生阻力损失的根本原因。

流体流经弯头、阀门、突然扩大/缩小等管件时，流速大小或流动方向发生改变，形成漩涡，造成机械能损失，即为局部阻力损失。

弯头、阀门、进口、出口等局部产生的

\begin{equation}
h_\mathrm{f'} = \zeta \frac{u^2}{2}
\end{equation}

$\zeta$ 为局部阻力系数，与管件结构形式有关。

\begin{equation}
h_\mathrm{f'} = \lambda \frac{\sigma l_e}{d} \frac{u^2}{2}
\label{eq:local-resistance}
\end{equation}

\chapter{总阻力损失的计算}

总阻力为直管阻力与局部阻力叠加：

\begin{equation}
h_{\mathrm{f},\mathrm{total}} = \sum \lambda\frac{l}{d}\frac{u^2}{2} + \sum \zeta\frac{u^2}{2}
\label{eq:total-resistance}
\end{equation}
或写为
\begin{equation}
h_{\mathrm{f}} = h_{\mathrm{f}} + h^\mathbf{'}_{\mathrm{f}} = \left(\lambda\frac{l + \sum l_\mathrm{e}{d}}{d} + \sum \zeta \right) \cdot \frac{u^2}{2}
\end{equation}

对直管阻力，如果管段内的流速不同，需分段计算。

对同一管件，阻力系数或当量长度二选一

\bluebold{减少阻力损失的方法}

管理尽可能短，尽量走直线，少走弯路

尽量不安装不必要的管件和阀门

采用大直径管

\end{document}
